\section{The 90-Minute Cycle}\label{sec:cycle}
Human sleep architecture is characterized by a cyclical alternation between two primary states: Non-Rapid Eye Movement (NREM) sleep and Rapid Eye Movement (REM) sleep, as detailed in Sections~\ref{sec:nrem} and~\ref{sec:rem}.

These NREM and REM states do not occur randomly. They are organized into cycles, each lasting approximately 90 to 110 minutes, repeating 4--6 times per night~\cite{fuller2006}.

A typical cycle progresses as follows:
\begin{enumerate}
    \item Descent from N1 $\rightarrow$ N2 $\rightarrow$ N3.
    \item Ascent back from N3 $\rightarrow$ N2.
    \item The first REM period, which is typically short (5--10 minutes).
\end{enumerate}
As the night progresses, the architecture of this cycle changes:
\begin{itemize}
    \item \textbf{NREM Dominance (Early Night):} The first few cycles are dominated by deep N3 (Slow-Wave) sleep, prioritizing physical restoration and declarative memory consolidation.
    \item \textbf{REM Dominance (Late Night):} In the latter half of the night, N3 sleep diminishes, and the REM periods become progressively longer and more intense, often lasting 30--60 minutes by the early morning. This prioritizes emotional processing and procedural skill generalization.
\end{itemize}

This reliable, two-part structure is not an accident. It is a fundamental algorithm for learning, which forms the core thesis of this paper.

\vfill
\textit{This chapter provided the ``what'' of sleep architecture. The next chapter explores the ``why,'' detailing the distinct functional roles of these two primary states.}
